% LAB BOREDOM — module L7, the Close. GUIDED DRAFTING (desktop recipe). Session 2026-08-21.
% Spec: outline v8 §L7. Plan approved 2026-08-21: one batch, ¶¶23–24; DEPLOYMENT SANCTIONED
%   in both ¶¶ (this module IS the conclusion of Part III per the 2026-07-29 arc ruling —
%   desktop first, laboratory second, coda here — so the author's bar on deployment
%   discussion "until the conclusion of Part III" lifts exactly here).
% Cautions binding on this batch (spec + standing rulings): no re-arguing · "demonstrated"
%   never "prove", to the last line · no echo of the desktop's closing construction ·
%   data-as-agent phrasing · never "laboratory" ("the experiment"/"the physiological
%   experiment") · no film comparisons · no "ones" · no withdrawal narratives.
% Drafting L7.2 UNBLOCKS three deferred desktop items (ledgered, post-section tasks, not
%   prose): M2.3's revision · the arc-close question · the M7.1-vs-III-5 three-state
%   harmonization.
%
% BATCH 10 = ¶¶23–24. Drafted 2026-08-21, awaiting author review.
% Verification ledger, batch 10:
%   ¶23 rests entirely on established, verified material: the doxa-gate apparatus (L4 ¶3),
%     the nine divergent episodes (keystone-decomposition), the deployment's single-channel
%     reading (desktop M4.1, carried at characterization altitude only), the D-33/0b
%     principle (quoted at L4 ¶10, paraphrased here).
%   ¶24 characterizations follow the arc ruling's own terms (vocabulary held / frame doing
%     work on evidence; divergence = demonstration, not measurement). No new figures.
% BATCH 10 REV 1 (2026-08-21, author notes):
%   P23-1 "the ratification this section has traced" -> named plainly: doxa's verdict.
%   P23-2 the channel term unpacked and the measures NAMED (eye closure and gaze deviation);
%     "channel" rendered as "kinds of evidence" (what a person says vs a recording of the
%     watching body) so the term cannot be misread as the individual measures.
%   P23-3/P24-1 "deployment" REPLACED with the specific name: the classroom study of the
%     virtual learning environment / the classroom study.
%   ★ P23-4 PURPOSE CORRECTION (author): the classroom study was never designed to analyze
%     boredom/engagement — it studied the VLE and students' experience of it; boredom in its
%     surveys is HAPPENSTANCE, arriving unbidden in free text. ¶23 rebuilt on that fact
%     (which sharpens the joining: unprompted culprit-naming IS the first form
%     self-announcing; the experiment then purpose-built the two-kind recording).
% ★ ¶¶23–24 APPROVED by author 2026-08-21 (rev 1). ★★ MODULE L7 COMPLETE — 2 ¶¶. The three
%   deferred desktop items are now UNBLOCKED (M2.3 revision · arc-close · M7.1-vs-III-5
%   three-state harmonization) — post-section tasks. NOT committed.

The two studies of this Part answer for each other. The classroom study of the virtual
learning environment was never designed to analyze boredom or engagement: its surveys asked
students about the environment and their experience of using it, and wherever boredom
appears in their answers it arrives unbidden---a student volunteering that the footage was
dull, that the world offered nothing to do, that YouTube was simpler. Evidence of that kind
is one channel only: what a person says about an episode after living it, \textit{doxa}'s
verdict with nothing recorded beside it. One channel can attest the first form of
boredom---the form whose culprit is nameable and whose affirmation is immediate, the form
students name unprompted---and it can do no more. The second form requires two kinds of
evidence that can disagree about one and the same episode: a recording of the watching body
while the episode runs---here, eye closure and gaze deviation---set against the verdict the
person delivers afterward. The physiological experiment was built to that shape, and in nine
of its thirty-six episodes the two kinds part company: the second form's structural grammar,
carried by measurements. The third form neither study reaches, and the reason is principled
rather than instrumental: an attunement cannot be ascertained and ought not to be, so no
added recording and no larger cohort would bring it into the data. Together the two studies
mark out how far instrumented evidence about boredom extends: further than the field has
taken it, and less far than the field would like.

What the two halves establish are different kinds of things, and the difference is the
finding this Part hands forward. The classroom study put the frame's vocabulary to work on the
sparest authored world in this dissertation, at the scale of a cohort, and the vocabulary
held: channels, worlds, the chain, the forms---each found its referent in the collected
responses without strain. The physiological experiment posed a different test: whether a
philosophical account can be brought to data collected under a different and partly
incompatible definition, and made to read out of that data what the collecting definition
could not name. The nine divergent episodes are the answer. The first study is a frame
surviving contact with evidence; the second is a frame doing work on evidence---and the work,
not the survival, is what this Part carries forward. That is also why the central result of
this section is a demonstration rather than a measurement: no instrument measured the second
form, and none could; what the data demonstrated is that where the eyes and the verdict part
company, the shape of their parting is the shape Heidegger describes.
